% Supplementary S5: placement grid, as printed (Amendments A1-A4). LaTeX table source; count macros are defined in the manuscript.
% Generated by tools/placement/build.py — do not edit by hand; edit the spec and re-run.
\begin{table*}[!htbp]
\centering
\caption{Placement of the content of Table~\ref{tab:competency_domains} across the three phases. Each entry gives a short treatment and, in brackets, the evidence grade of that placement under the rules of Section~\ref{sec:methods_ages}: H, learners in the phase's age band (reported, or inferred from the grade and the named country) were taught and tested on the element, its prerequisites, if any, appear in standard curricula by that phase, and no consensus source opposes it; M, evidence only from an adjacent age band, or an opposing consensus; L, contrary evidence; Conj, no evidence from the phase or an adjacent band, so the placement is a conjecture to be tested. Ages are indicative. Of \nFtEvidence{} eligible studies with outcomes, \nFtTaughtTested{} taught and tested at least one element; combining reports of the same sample leaves \nPlSamples{} samples. The full matrix (\nPlCells{} placed cells), with sources and page numbers for every cell, is in Supplementary S5. $^{a}$Conditional domains.}
\label{tab:framework_pathway_map}
\scriptsize
\begin{tabularx}{\textwidth}{@{}>{\raggedright\arraybackslash}p{2.3cm}>{\raggedright\arraybackslash}X>{\raggedright\arraybackslash}X>{\raggedright\arraybackslash}X@{}}
\toprule
Domain & Phase~1 (indicative ages 11--13) & Phase~2 (14--16) & Phase~3 (16--18) \\
\midrule
\rowcolor{dphen!30} D1 Principles and phenomena & superposition (verbal) [M]; single-photon MZI, qualitative [M]; random single outcomes [H]; uncertainty (verbal) [M]; wave and particle behaviour [M]; correlated pairs (verbal) [M]; no path, no prediction [M] & superposition, real two-path states [H]; single-photon MZI, probabilities [H]; outcome frequencies [H]; single-slit trade-off [H]; photons build interference [H]; pair-outcome probabilities [H]; no-trajectory argument [H] & superposition, real-valued states [H]; MZI, real amplitudes (0/$\pi$) [H]; probabilistic prediction [H]; Heisenberg estimate [H]; de Broglie wavelength [H]; entangled states, real [H]; state vs property [H] \\
\rowcolor{dmath!30} D2 Mathematical representations & -- & probabilities as frequencies [H]; real two-state vectors [H]; particle in a box (qualitative) [M]; Dirac notation, real coefficients [M] & probabilities of sequences [H]; real vectors in two bases [H]; particle in a box (qualitative) [M]; complex amplitudes (optional) [Conj]; Dirac notation, real coefficients [H]; $2\times2$ real matrices (optional) [M] \\
\rowcolor{datom!30} D3 Atoms, light, and matter & line spectra seen [M]; photons (verbal) [H]; emission, absorption (verbal) [Conj]; semiconductors named [Conj] & level diagrams, $E=hf$ [H]; $E=hf$, photoelectric effect [H]; transitions between levels [M]; energy-gap picture [M] & hydrogen levels, spectra [H]; photoelectric equation [H]; stimulated emission, laser [H]; bands from levels [H] \\
\rowcolor{dinfo!30} D4 Information and qubits & bit vs qubit (descriptive) [M]; qubit read-out 0/1 [M]; entanglement as resource (named) [M]; qubit platforms named [M] & qubit as real vector [H]; read-out probabilities [H]; shared key from pairs [H]; what encodes 0 and 1 [H] & states of $n$ qubits [H]; read-out in two bases [H]; teleportation (conceptual) [M]; isolation and cooling [H] \\
\rowcolor{dcomp!30} D5 Computing and simulation & what QCs can and cannot do [M]; visual circuit blocks [M]; fragility, decoherence named [M] & which problems gain [H]; X, H, CNOT in a simulator [H]; errors as flipped outcomes [H] & speed-up claims assessed [M]; small circuits (conceptual) [H]; error-correction idea [M] \\
\rowcolor{dcomm!30} D6 Communication$^{a}$ & eavesdropper leaves traces [Conj]; PQC and QKD named [Conj]; quantum randomness [H] & wrong setting randomises [M]; BB84 walkthrough [M]; how PQC and QKD differ [Conj]; beam-splitter coin [H] & eavesdropper detection [H]; BB84 walkthrough [H]; PQC and QKD compared [Conj]; how a QRNG works [M] \\
\rowcolor{dsens!30} D7 Sensing and metrology$^{a}$ & clocks for navigation [Conj]; sensitivity makes sensors [Conj]; sensors and uses named [Conj] & clock from an atomic transition [Conj]; phase sensing with an MZI [M]; what sensors measure [M]; atom interferometry (awareness) [Conj] & clocks from energy levels [Conj]; interference and spin sensing [M]; quantum vs classical sensors [M]; atom interferometry (awareness) [Conj] \\
\rowcolor{ddev!30} D8 Devices and enabling technologies & quantum devices named [M]; photon detectors, scanners [Conj]; cold, vacuum, sources named [Conj] & LEDs, solar cells, $E=hf$ [H]; photodetector principle [M]; why cooling isolates [Conj] & laser principle, bands [M]; coincidence detection, MRI spin [H]; photon sources, cryogenics [Conj] \\
\rowcolor{dnos!30} D9 Nature, history, philosophy & discovery stories [Conj]; disagreement named [Conj]; photon as a model [Conj]; science changes [M] & changing models [Conj]; Schr\"odinger's cat [Conj]; limits of models [Conj]; determinism revisited [H] & Bohr--Einstein debate [Conj]; interpretations compared [Conj]; models vs reality [Conj]; paradigm change [M] \\
Practice layer & LEDs, filters, diffraction [Conj]; observe simulations [Conj]; block-based circuits [Conj] & Planck's constant with LEDs [Conj]; test predictions [M]; program small circuits [Conj] & experiments with uncertainty [Conj]; ``what if'' simulations [H]; high-level toolkit [Conj] \\
Societal dimension & why encryption must change [Conj]; sensing benefits named [Conj]; ``quantum'' in adverts [Conj] & PQC migration [Conj]; benefits and limits [Conj]; checking health claims [Conj] & migration trade-offs [Conj]; who benefits [Conj]; pseudoscience analysed [Conj] \\
\bottomrule
\end{tabularx}
\end{table*}
